\DocumentMetadata{
  pdfversion=2.0,
  pdfstandard=ua-2,
  lang=en-US,
  tagging=on,
  tagging-setup={math/setup=mathml-SE,tabsorder=structure}
}
\documentclass[11pt,letterpaper]{article}
\usepackage[margin=0.68in,includefoot]{geometry}
\usepackage{newcomputermodern}
\usepackage{amsmath,amscd,mathtools}
\usepackage{tikz,adjustbox}
\providecommand{\square}{\mdlgwhtsquare}
\usepackage[hidelinks]{hyperref}
\hypersetup{
  pdftitle={Analysis PhD Exam, August 2023},
  pdfauthor={University of Florida Department of Mathematics},
  pdfsubject={Graduate mathematics examination},
  pdfdisplaydoctitle=true,
  bookmarksopen=true
}
\setcounter{secnumdepth}{0}
\setlength{\parindent}{0pt}
\setlength{\parskip}{0.28em}
\setlength{\emergencystretch}{3em}
\setlength{\leftmargini}{2.15em}
\setlength{\leftmarginii}{2.65em}
\setlength{\leftmarginiii}{3em}
\pagestyle{empty}
\raggedbottom
\allowdisplaybreaks
\NewTaggingSocketPlug{block/list/label}{examproblem}{%
  \tagstructbegin{tag=Lbl}\tagstructend%
  \tagstructbegin{tag=\LBody}%
  \tagstructbegin{tag=\ProblemHeadingTag,title-o={\ProblemHeadingText}}%
  \tagmcbegin{tag=\ProblemHeadingTag,actualtext={\ProblemHeadingText}}%
  #2%
  \tagmcend%
  \tagstructend%
}
\newenvironment{examproblems}[2]{%
  \def\ProblemHeadingTag{#2}%
  \AssignTaggingSocketPlug{block/list/label}{examproblem}%
  \begin{enumerate}%
  \setcounter{enumi}{#1}%
  \renewcommand{\labelenumi}{\textbf{\theenumi.}}%
  \setlength{\topsep}{0.35em}%
  \setlength{\partopsep}{0pt}%
  \setlength{\itemsep}{0.48em}%
  \setlength{\parsep}{0.18em}%
}{\end{enumerate}}
\newenvironment{examsubparts}{%
  \AssignTaggingSocketPlug{block/list/label}{default}%
  \begin{enumerate}%
  \setlength{\topsep}{0.18em}%
  \setlength{\partopsep}{0pt}%
  \setlength{\itemsep}{0.16em}%
  \setlength{\parsep}{0.1em}%
}{\end{enumerate}}
\newenvironment{examinstructions}{%
  \par\begingroup\itshape
}{%
  \par\endgroup\vspace{0.18em}
}
\makeatletter
\renewcommand\subsection{\@startsection{subsection}{2}{0pt}%
  {-1.25ex plus -.25ex minus -.1ex}%
  {0.55ex plus .1ex}%
  {\normalfont\large\bfseries}}
\renewcommand{\@maketitle}{%
  \newpage\null\vskip -1em%
  {\footnotesize\bfseries
    \href{https://www.ufl.edu/}{University of Florida}, \href{https://math.ufl.edu/}{Department of Mathematics}\par}%
  \vskip -0.5em%
  \tagstructbegin{tag=H1,title={Analysis PhD Exam, August 2023}}%
  \tagpdfparaOff%
  \tagmcbegin{tag=H1}%
    {\LARGE\bfseries\href{https://gma.math.ufl.edu/exams/analysis-phd/}{Analysis PhD Exam}, August 2023\par}%
  \tagmcend%
  \tagpdfparaOn%
  \tagstructend%
  \vskip 0.25em%
  \tagmcbegin{artifact}%
    \hrule height 1pt%
  \tagmcend%
  \vskip 1em%
}
\makeatother
\title{Analysis PhD Exam, August 2023}
\author{}
\date{}
\begin{document}
\maketitle
\thispagestyle{empty}
\begin{examinstructions}
Write solutions in a neat and logical fashion, giving complete reasons for all steps. You must attempt SIX problems.
\end{examinstructions}
\begin{examproblems}{0}{H2}
\def\ProblemHeadingText{Problem 1.}
\item
State Tonelli’s theorem, and give an example to show that the~\(\sigma\)-finiteness hypothesis is necessary.
\def\ProblemHeadingText{Problem 2.}
\item
Give an example of each of the following, if possible. (If not possible, give a brief explanation of why.) When giving examples, be clear about what measure space you are working in.
\begin{examsubparts}
\item[\textnormal{a)}]
a sequence of functions converging in measure but not in \(L^1\),
\item[\textnormal{b)}]
a sequence converging almost uniformly but not essentially uniformly,
\item[\textnormal{c)}]
a sequence converging in \(L^1\) but not almost everywhere.
\end{examsubparts}
\def\ProblemHeadingText{Problem 3.}
\item
Let \((X,\mathcal{M},\mu)\) be a measure space. Prove that \(L^2(\mu)\) is complete (in the \(L^2\) norm).
\def\ProblemHeadingText{Problem 4.}
\item
Let~\(\mathcal{X}\) be a Banach space. Prove that every finite-dimensional subspace \(\mathcal{M}\subset\mathcal{X}\) is closed.
\def\ProblemHeadingText{Problem 5.}
\item
This problem has two parts.
\begin{examsubparts}
\item[\textnormal{a)}]
State the Lebesgue–Radon–Nikodym theorem.
\item[\textnormal{b)}]
State and prove a version of the chain rule for Radon–Nikodym derivatives.
\end{examsubparts}
\def\ProblemHeadingText{Problem 6.}
\item
Let \((X,\mathcal{M},\mu)\) be a~\(\sigma\)-finite measure space, let \(1\leq q<p<\infty\) and suppose that \(L^q(\mu)\subseteq L^p(\mu)\). Prove that there exists a constant \(c>0\) such that, for every \(E\in\mathcal{M}\), either \(\mu(E)=0\) or \(\mu(E)\geq c\).
\def\ProblemHeadingText{Problem 7.}
\item
Let~\(\mathcal{H}\) be a Hilbert space and let \(A:\mathcal{H}\to\mathcal{H}\) be a linear map (not assumed bounded). Suppose there exists another linear map \(B:\mathcal{H}\to\mathcal{H}\) (not assumed bounded) so that \(\langle Ax,y\rangle=\langle x,By\rangle\) for all \(x,y\in\mathcal{H}\). Prove that~\(A\) and~\(B\) are bounded.
\def\ProblemHeadingText{Problem 8.}
\item
This problem has two parts.
\begin{examsubparts}
\item[\textnormal{a)}]
State the Closed Graph Theorem and Banach Isomorphism theorem.
\item[\textnormal{b)}]
Using the Banach Isomorphism theorem (or otherwise), prove the Closed Graph Theorem.
\end{examsubparts}
\end{examproblems}
\end{document}
